\section{Technical Implementation}

\bigskip
\bigskip

This section describes how we chose to develop our app with a special focus on the technical parts or the strong choices made. 

\bigskip
\bigskip

\subsection{Technical Stack and Development Environment}

\bigskip
\bigskip

The application is coded mainly in Python 3.13 and some parts are in SQL which are the ones about data, like searching directly in the database or modifying it.

\bigskip

The biggest part of the backend is divided into three layers: the DAO, services, controllers. For each function, a DAO, service and controller method is developed. The Data Access Object methods are the only ones that communicate with the database by using SQL queries. Then what is sent back from the DAO goes to the service which processes this information, in order to give a clear response to the controller and then to the user through the frontend.

\medskip

One major part of the backend is about FastAPI. It is a tool used to create web API easily. The first role of it is to send the request wanted by the user to the app and send the response back to the user in JSON.

\bigskip

In the controller folder of the backend, each function in all the files becomes a URL that the frontend (Streamlit) can call. For example, when somebody clicks on the "Login" button, Streamlit sends an HTTP request to the \texttt{/login} address generated by FastAPI. That makes the code more readable by dividing each function (or group of functions) in a precise file, almost totally independent from other parts. 

\bigskip

Furthermore, to store our data, we use PostgreSQL on Onyxia in a dedicated schema. It is where the DAO sends their SQL requests. The connection between PostgreSQL and VSCode is made through  a \texttt{.env} file, with all the needed connection information listed in the \texttt{README}. Then, not for users but for us during the development, we use CloudBeaver. That is an interface to have a look at the tables. 

\bigskip

When we first launch the app, of course there is nothing in the database. At startup, a FastAPI \texttt{lifespan} handler runs \texttt{run\_init\_db()} which initialises all the necessary tables and a default admin, or at least verifies that they already exist.



\subsection{NASA API Integration and Data Synchronization}

\bigskip
\bigskip

In order to use our app we obviously need some data. We import them from the API NeoWs of NASA. It records the approaches of NEOs around the Earth and all the information linked with the object, like the distance, the speed or the size.

\bigskip

In this API, there is an endpoint \texttt{/feed} that displays the NEOs detected on a period of maximum seven days. That is why, in \texttt{import\_service}, the method \texttt{import\_neo\_range} automatically splits any requested period into consecutive seven-day chunks. Moreover, there is a second method named \texttt{import\_all} that uses the first one by giving it a period of three years : two years before today and one after. That represents about 160 API calls that we have spaced a bit to avoid any restriction of NeoWs. Additionally, the imports run in a background thread. The background starts an import when the web service starts, then with the \texttt{DailyImportScheduler}, every 24 hours there is an update or if an administrator clicks on the button. During this time, NEOs appear little by little in app.

\bigskip

Once these NEOs arrive in our app, how do we store them ? First we complete the table \texttt{neo} with some information like their size as mentioned earlier or if they are considered as hazardous by NASA. However, as the responses of NeoWs are not pure NEOs but a list of approaches, we also complete the table \texttt{approach} where we can find all the approaches. And one NEO can have multiple approaches recorded which will help us with some functions later. But obviously each \texttt{neo} and \texttt{approach} are unique by using the SQL command \texttt{ON CONFLICT}.


In particular, the object created by users is of course saved in the table \texttt{neo} with a label 'USER'.
    
\bigskip

Finally, for our function APOD (Astronomy Picture of the Day) we do not store pictures across time, we allow users to see the APOD through a button going to the NASA site. We did it like this because currently, the API returns just the NASA logo instead of the APOD. Users can still read the description and then go to the website.

\bigskip
\bigskip

\subsection{Security and Access Control}

\bigskip
\bigskip

In this part, we will present some security actions that we have used for the app.

\bigskip

The first one is of course the password hashing. For some obvious reasons, we do not store plaintext passwords. In fact, we store the result of \texttt{hash\_password}. It computes a character string by a salted SHA-256 algorithm. We can store this result as it is not reversible even if a person knows the encryption algorithm. Although this algorithm is fast, it remains weak and vulnerable. Some other methods could have been better like \texttt{bcrypt} or \texttt{argon2}.

\bigskip

Then, once the password and username have been checked, the backend generates a JSON Web Token (JWT) to authenticate next requests without recalling username and password. There are two tokens, first an access token, available for 30 minutes, and a refresh token for seven days automatically refreshed. And to follow the connection, we have access to the list of all the previous logins through the admin menu in app using \texttt{ConnectionService} and \texttt{ConnectionDao}.

\bigskip

Furthermore, some functions are restricted, only an administrator has access to them. For example, we can reset the database in app, but this action has huge consequences, that is why its access is limited. We have the same principle with \texttt{check\_same\_user} and \texttt{\_check\_alert\_owner} which prevent users from modifying other users' alerts or favourites. Otherwise, an error \texttt{403 Forbidden} is returned.

\bigskip

Finally, the SQL injection is prevented in all the DAO classes like \texttt{NeoDao}, \texttt{ApproachDao}, \texttt{UserDao} or \texttt{AlertDao}. They rely exclusively on parameterized SQL queries with the module \texttt{psycopg2}.

\subsection{Core Business Logic Implementation}

\bigskip
\bigskip

In this part, we will see some interesting functions that we have done, which are worth describing.

\bigskip

Once our tables are full, we do not call NeoWs anymore. For example, in \texttt{NeoService} our method of Multi-criteria search only searches in our database. In order that the user finds NEOs according to his criteria, the function can compose filters in the SQL query. This filter can be composed of many categories that we find in \texttt{criteria\_map} in \texttt{NeoDao}. We have for instance \texttt{neo\_id}, the name, the size or even the approach date.

\bigskip

Also, for \texttt{StatisticsService} and \texttt{DashboardService}, to avoid importing all we have in the table, and then processing them in the service, our SQL queries already filter the data we need and sometimes, if the command is simple enough, process some information. For example, when we need to count a number of elements, we could have returned to the service all the rows and count them, but we just use \texttt{COUNT} in SQL. However, when we do not have to use a lot of data for a method, we have tried to keep all the treatment inside the service.

\bigskip
\bigskip

Now, let's focus on the \texttt{Favourite} feature. It works by associating a \texttt{neo\_id} with a \texttt{user\_id}. Of course users can add or remove them from their favourite list. But users also need to display a graph. It has to show approaches of a NEO across time. That is exactly the role of our table \texttt{approach}. As in this table we already have all the approaches, their date, the distance from Earth and of course the \texttt{neo\_id}, it is quite easy to plot that kind of graph. First, we have to get all the approaches of our favourite. That is why we have a function \texttt{approaches\_favourite} which returns all approaches of favourites given a user in a dictionary (with \texttt{neo\_id} as a key associated with a list of \texttt{Approach}). Then in \texttt{graph\_favourite}, we use that return, only where the key is the NEO selected. Since there is now only a list of \texttt{Approach}, we plot the graph of the distance of the NEO according to the time. And to send the graph to the frontend, we convert the \texttt{png} file into a base64 string as JSON only carries characters or numbers. And the streamlit frontend decodes it to be able to display the graph to the user.

Unfortunately, as we import data within a period of three years, many NEOs have a periodicity above three years, which led us to graphs with only one point. Hopefully some are more interesting, with many approaches during these three years, and consequently a more informative plot.

\bigskip

Finally, this section will describe our feature about alerts. One decision that we made, is to consider the alerts only on the NEOs in the table \texttt{approach}. We consider that the only ones that could be worth to be notified about are the ones that could approach in the future. Therefore, we will consider only approaches after the current date. To set an alert, users can give values about a criterion. These criteria are: \texttt{min\_diameter}, \texttt{max\_distance} and \texttt{hazardous\_only}. The first one filters the NEOs by their diameter if it is above the value, same with the distance if the distance is below the value, and \texttt{hazardous\_only} filters by the hazardous ones (if no value, it is \texttt{False}). Of course, to create an alert, it is mandatory to give at least one criterion, otherwise, all the upcoming \texttt{approach} would be considered in the alert.

Once it is done we store it in the table \texttt{alert}. We filter the matching approaches directly in \texttt{AlertDao} and return them with the method \texttt{check\_send\_notification} which returns a notification that has been raised (or nothing if no matches). Furthermore, to avoid the user being overwhelmed with notifications, each time it is triggered, the \texttt{last\_eval\_date} updates itself, and cannot be triggered again until the next session. And even if the user is tired of one notification that keeps being triggered every day, he can turn it off without deleting it. When he wants to have new notifications with the same criteria, he can now simply turn it on.


\subsection{What we planned and what we did}

In September, in the first analysis report we fixed the main technical choices. It included twelve endpoints all gathered under a common prefix \texttt{/api/v1}. All personal data should have been accessible through a path like \texttt{/users/me}. Also, we planned that the NASA data would be loaded by a daily batch that adds one more day after the imported period.

\medskip

However, even if the global architecture did not change, these details did. In fact, now, there are thirty endpoints that you can find in the Table of endpoints (Table 4).

Furthermore, in the final version we use the identifier of the user in the path about personal data, for example we have \texttt{/users/12} instead of \texttt{/users/me}. Leaving the situation like that would have allowed other users to access personal data of everybody just by modifying the id in the endpoint. That is why the server compares this id with the one of the token. Administrators are the exception, actually they can access and manage every account.

Some other noticeable changes, the daily update do not only work every day. It can update the database starting at the last import done, restart at the right point if the last import did not finish. We can see that when the status of the import is \texttt{partial} instead of \texttt{success} or \texttt{failed}. And also reloads the next thirty days in case that new NEOs have been found in the meantime. Lastly, an import is not anonymous anymore as the table \texttt{import} requires an id, by default it is the admin (in case of an automatic update).

\bigskip
\bigskip